%  Contents/dados_no_overgrown.tex
\midsectiontitle{Dados no Overgrown}

\section*{A Manifestação Numérica}
No universo de \textit{Overgrown}, os dados transcendem a mera casualidade, eles representam a manifestação das capacidades inatas e divinas de um personagem. Cada rolagem é o encontro entre o potencial do aventureiro, os conhecimentos do Pináculo e a pressão do momento presente.

\section*{Os Três Pilares das Rodadas}
Toda interação significativa com o mundo é regida por uma estrutura de três fundamentos que ditam o destino de uma jogada:

\begin{description}
    \item[\ringm{Potencial Inato (Atributos)}:] Refletem a força bruta, agilidade ou a sabedoria divina. No sistema \textit{Overgrown}, o valor do seu atributo não é apenas um bônus, mas o que determina o volume de poder, ou seja, a quantidade de dados, que você pode manifestar em uma ação.
    
    \item[\ringm{Modificadores de Fluxo}:] São as variáveis externas. Bônus provenientes de equipamentos, bênçãos de magias, trilhas da Árvore de Talentos ou vantagens táticas situacionais concedidas pelo mestre.
    
    \item[\ringm{A Resultante Final}:] O ápice da jogada. É a soma do melhor resultado obtido nos dados com os modificadores, definindo se as ações do personagem se empoem sobre a realidade ou se foi superada pelos desafios encontrados.
\end{description}

\section*{Mecânica de Resolução}
O sistema utiliza principalmente o dado de vinte lados (\textbf{d20}) para testes de perícia, ataques e salvaguardas. Dados de outras faces, do \textbf{d2} (Moeda) ao \textbf{d100}, são reservados para calcular o impacto de danos, efeitos de magias e eventos aleatórios do destino.

\subsection*{Ascensão de Atributos}
A maestria em um atributo concede ao personagem uma vantagem grandiosa. A cada \textbf{5 pontos} investidos em um atributo, o aventureiro ganha o direito de arremessar um dado adicional, escolhendo o maior resultado entre eles. Esta progressão não apenas eleva a média dos resultados, mas expande exponencialmente as chances de manifestar um Sucesso Crítico.

\begin{rpgboxtips}{Intervenção Crítica}
Um resultado \textbf{20} natural é um Sucesso Crítico: um acerto garantido que ignora qualquer tentativa de esquiva, podendo ser mitigado apenas por bloqueios. Por outro lado, um \textbf{1} natural é uma Falha Crítica, um erro catastrófico que rompe o fluxo da ação. Lembre-se: cada ataque e magia possui efeitos críticos únicos em suas descrições.
\end{rpgboxtips}

% Tabela de Progressão de Poder
\begin{center}
\renewcommand{\arraystretch}{1.3}
\begin{tabular}{@{} cc @{}}
\toprule
\rowcolor{gray!15} \textbf{\ringm{Magnitude do Atributo}} & \textbf{\ringm{Reserva de Dados (d20)}} \\ 
\midrule
1 -- 5   & 1d20 \\
6 -- 10   & 2d20 \\
11 -- 15 & 3d20 \\
15 -- 20 & 4d20 \\
\textbf{21+} & \textbf{5d20} \\ 
\bottomrule
\end{tabular}
\end{center}

\begin{rpgboxtips}{Além do Limite Mortal}
O limite natural de um aventureiro é a reserva de \textbf{5 dados}. Contudo, aqueles que portam Selos de Divindade ou acessam estados de superiores através de \textbf{Nós Supremos} podem quebrar esta barreira, manifestando reservas de poder que desafiam a compreensão humana. Ter 0 atributos resulta em desvantagem rolando dois dados e pegando o menos.
\end{rpgboxtips}

\begin{comment}
    

\begin{armaCard}{ESPADA DO PINÁCULO}
    \textit{\small Uma lâmina forjada em Verdun, banhada em energia primordial que ressoa com a vontade do portador.}
    
    \vspace{8pt}
    
    % Linha 1 de Atributos
    \attrField{Base Damage}{2D8}
    \attrField{Attribute}{FOR}
    
    \vspace{2pt}
    
    % Linha 2 de Atributos
    \attrField{Multiplier}{2*}
    \attrField{Total DMG}{2D8 + 2*FOR}
    
    \vspace{8pt}
    \overRule
    \vspace{4pt}
    
    % Bloco de MIG, GRA, FOR, WIS, SEN
    \statBox{MIG}{12}
    \statBox{GRA}{08}
    \statBox{FOR}{10}
    \statBox{WIS}{05}
    \statBox{SEN}{06}
    
    \vspace{8pt}
    \overRule
    \vspace{4pt}
    
    % Efeitos Mágicos
    \textbf{\textit{Lâmina de Éter.}} O dano desta arma ignora resistências físicas de seres não-divinos.
    
    \vspace{10pt}
    {\ringmfamily\Large\color{overgrownRed} ACTIONS}
    \overRule
    
    \textbf{\textit{Corte Dimensional (2 Ações).}} O jogador realiza um ataque em área (cone de 3m). Inimigos atingidos devem realizar um teste de \textbf{GRA} ou sofrerão o dano total da arma.
    
\end{armaCard}
\end{comment}